\chapter{Summary and Conclusions}
\label{ch:conclusions}

Spain has more than 34.7 million registered vehicles and one of the highest vehicle densities in the European Union, and the problem this creates in urban areas is less about an absolute shortage of parking spaces than about a mismatch in their distribution, since public streets are saturated while private garages built under residential and commercial buildings sit empty for most of the working day. In the Canary Islands the situation is compounded by insularity, which limits viable public transport alternatives and makes the private car the dominant mode of urban mobility. Tenerife registers 870 vehicles per thousand inhabitants, the highest rate in the archipelago, and both Santa Cruz de Tenerife and San Cristóbal de La Laguna concentrate most of their commercial, university and administrative activity in historic districts whose street layout predates mass car ownership, leaving very limited room for off-street public parking supply. The ZBE ordinance approved in April 2026, which restricts access to vehicles without a DGT environmental label, adds further pressure by pushing demand toward private spaces in the urban perimeter, a segment that no existing platform is currently structured to serve.

This work aimed to demonstrate, on one hand, that the complete transaction cycle of a peer-to-peer parking marketplace can be resolved within a single application without requiring any coordination between parties outside it, covering space discovery, booking with database-level conflict validation, payment with automatic IGIC calculation and physical garage access, and on the other that doing so produces a commercially viable product within a realistic investment horizon. A subsidiary question was whether the fiscal and regulatory specificity of the Canarian market, in particular the IGIC regime and the ZBE compliance logic, could be handled natively in the architecture rather than worked around.

The system was designed so that the invariants that matter most for correctness are enforced at the lowest possible layer. Booking overlap validation runs as a PL/pgSQL trigger within the database transaction, so conflicting reservations are rejected atomically regardless of the number of concurrent requests, rather than relying on application-level checks that could be bypassed under load. Row Level Security policies in PostgreSQL enforce tenant isolation at the database engine level, independently of the service layer, so a bug in application code cannot expose one user's data to another, while the IGIC calculation and the Stripe Connect payment split are computed in Supabase Edge Functions on the server side, which prevents client-side manipulation of the fiscal breakdown and keeps business logic out of the frontend. Physical garage access is handled through SmartAccess, which generates time-bounded tokens linked to the active booking state and can verify them without a persistent network connection, and the entire product, web application and Android APK alike, is built from a single TypeScript codebase via Capacitor, with React~18 concurrent mode keeping the Leaflet map responsive when rendering large numbers of markers simultaneously.

The commercial study puts the total investment required to bring the prototype to a market-ready product at 287,000\,€, covering a professional team across analysis, design, development, testing and deployment. Revenue comes from an asymmetric transaction commission on each booking, with 15\,\% charged to the driver and 8\,\% to the owner, which yields a net platform income of approximately 1.77\,€ per booking after Stripe fees, together with a Parky Pro subscription for multi-space landlords at a weighted average of 20.99\,€ per month and premium visibility micropayments averaging 3.19\,€ per listed space. Under a conservative growth scenario of 100 active spaces at launch growing to approximately 990 by month twenty, with demand stable at 15 bookings per space per month, the model reaches operational break-even in month three and recovers the full initial investment in month nineteen, with an accumulated surplus of 157,038\,€ at the end of month twenty-four, corresponding to a two-year ROI of 54.7\,\%.

The results of this work show that a full-featured peer-to-peer parking marketplace is technically feasible within a realistic budget, given that all seven specific objectives defined in Chapter~1 are met by the deployed prototype and that the architectural choices that push correctness invariants to the database layer reduce application complexity rather than just moving it elsewhere. The Canarian fiscal and regulatory specificity that might deter external competitors, most notably the native IGIC handling and the ZBE compliance logic, turns out to be an implementable advantage that raises a meaningful barrier for any platform entering from the mainland. Because the same infrastructure, tax structure and regulatory framework apply across the entire archipelago, geographic expansion does not require rebuilding the product, only replicating the market entry process in cities where the parking density problem is comparable.

% ============================================================
\section{Future work}
\label{sec:future_work}
% ============================================================

The most direct growth path is expansion across the Canary Islands archipelago, which requires no architectural changes since the technical infrastructure, business model and regulatory framework are identical in Las Palmas de Gran Canaria, Arrecife and every other island capital. Within the product itself, the extensions with the greatest impact on both revenue and user retention are those described in Chapter~\ref{sec:funcionalidades_adicionales}, including passive Bluetooth-based occupancy detection to automate availability updates, event passes and recurring subscriptions via Stripe Billing, automatic ZBE exemption validation, real-time demand-based pricing and a bidirectional reputation system that would address the information asymmetry inherent to any two-sided marketplace in its early stages.
